\documentclass[11pt]{article}
% Shared preamble for the round 1 response letters.
% RED-D-26-00078, submitted as "Housing First or Treatment First?",
% resubmitted as "Housing First or Treatment First? Ask the Kids, Not the Parents"
%
% House style, decided once here:
%   * responses are in RED, referee text is in the black italic block above them
%   * NO HYPHENS anywhere in the prose
%   * every response ends with a pointer to where in the new paper it lives

\usepackage[margin=1.05in]{geometry}
\usepackage{amsmath,amssymb}
\usepackage{xcolor}
\usepackage[official]{eurosym}
\usepackage{booktabs}
\usepackage{longtable}
\usepackage{array}
\usepackage[hidelinks]{hyperref}

% enumitem is not in this TeX installation, so first level list spacing is set
% directly instead
\makeatletter
\renewcommand{\@listI}{%
  \leftmargin=1.7em \labelwidth=1.2em
  \topsep=0.4em \parsep=0pt \itemsep=0.3em \partopsep=0pt}
\makeatother

\definecolor{respred}{RGB}{190,25,25}

\setlength{\parskip}{0.35em}
\setlength{\parindent}{0pt}

\newcolumntype{L}[1]{>{\raggedright\arraybackslash}p{#1}}
\newcommand{\pname}[1]{\texttt{#1}}

% ---- the two building blocks -------------------------------------------------
% \cmt{label}{short title} opens a comment, then the verbatim or paraphrased
% referee text goes in the referee environment, then the response.

\newcommand{\cmt}[2]{%
  \par\medskip
  \noindent\textbf{#1\quad #2}\par}

\newenvironment{referee}
  {\begin{list}{}{\leftmargin=1.6em \rightmargin=1.2em \topsep=0.4em}\item[]\small\itshape}
  {\end{list}}

\newenvironment{response}
  {\par\begingroup\color{respred}\medskip}
  {\par\endgroup\medskip}

% a flag for statements in the reports that are factually wrong, so they are
% easy to find again and easy to handle carefully
\newcommand{\correction}{\textbf{[Point of correction.]}\ }
\newcommand{\concede}{\textbf{[Conceded.]}\ }
\newcommand{\newresult}{\textbf{[New result.]}\ }
\newcommand{\open}{\textbf{[Still open.]}\ }

% numbers that are not in yet, so they are impossible to miss on a read through
\newcommand{\tbd}[1]{\textbf{[TO FILL: #1]}}

% Marks what has been done since the first draft of these letters. The letters
% were written when most of this was a plan; this flag marks the places where it
% is now built and running.
\newcommand{\thisround}{\textbf{[Built since the first draft of this letter.]}\ }

\newcommand{\lettertitle}[2]{%
  \begin{center}
  {\Large\bfseries Housing First or Treatment First?}\\[0.3em]
  {\normalsize Manuscript RED-D-26-00078}\\[0.9em]
  {\large\bfseries #1}\\[0.3em]
  {\normalsize #2}
  \end{center}
  \vspace{0.6em}}


\begin{document}

\lettertitle{Response to Referee 1}{Comments in order, responses in red}

Thank you for a detailed and careful report. You ordered your comments by
importance as 2, 1, 4, 3. I answer them below in the order of the report, and I
start with a short overview of what changed in the paper that bears on them.


{\color{respred}
\textbf{One note before the comments.} The paper has a new title: it was submitted as
\emph{Housing First or Treatment First?} and is now \emph{Housing First or Treatment First? Ask the Kids, Not the Parents}. The reason is the substance of the revision rather than a relabelling, and it
should become clear over the course of this letter.

\textbf{What changed that bears on your report.} The version you read ranked
rationed and universal programmes, and it put universal Treatment First ahead of
universal Housing First only once a loss of consortium was added to welfare after
the model was solved. Your comments 3(a) to 3(c) were right about that
adjustment. It is gone, and so are the universal programmes. The new paper
compares the two policies within the existing rationed programme, and it asks
first whether moral hazard, the argument that dominates the policy debate,
justifies a sobriety condition. It does not: with the externalities on children
switched off, a newborn is worse off under Treatment First by 1.09 percent of
lifetime consumption, even when the savings from the drug test are spent on more
vouchers. The answer turns instead on the children of users, through three externalities
that a parent counts only in part. A parent's drug spending, and the unemployment
that a substance use disorder (SUD) brings, take consumption and housing from the
child; a parent's SUD lowers the child's productivity as an adult; and a revoked
voucher can put the child in a shelter. With these in the model, and with the drug
price clearing its market and the chance of drawing a voucher adjusting to balance
the budget, Treatment First cuts drug use by 21 percent and homelessness by 11
percent, and raises a newborn's welfare by 2.57 percent of lifetime consumption,
three quarters of it an externality its parent does not count.

Many of your other comments are now answered in the model or in the calibration.
The quantitative model reports its fit on every targeted moment (Tables 3 and A1)
and its untargeted moments beside the data (Table 4). It now has serious mental
illness (SMI), a treatment choice, a drug price that clears its market, a drug
test that catches only some users and whose cost is paid from the budget, and a
transition that is solved and reported (Sections 6.4 and 6.5). A two period
stylized model with saving (Section 3) replaces the old static appendix, and the
comparison you ask for in 1(d) is run in it. I also removed endogenous labour
supply, rents, homeownership and house prices, which in earlier versions had a
negligible effect on the results (Section 4.4, footnote 19); this bears on 2(c)
and 2(d). Where the paper does not do what you asked, the response says so and is
marked as still open.
}

%======================================================================
\section*{1. Model}

\cmt{1(a)}{Mental health is largely absent}
\begin{referee}
``Mental health is central to the economic and policy environment being studied,
but it is largely absent from the model. \dots It seems to me that mental health
contributing to drug use and addiction as well as homelessness would weaken TF
dramatically as a policy.''
\end{referee}
\begin{response}
Added. SMI, mental illness severe enough to cause substantial functional
impairment, is now a state of the model, a two state chain. Onset follows an age
profile that reproduces a log linear fit to the share of adults with SMI at each
age in the 2024 National Survey on Drug Use and Health (NSDUH), scaled by a level calibrated to
its prevalence, 4.82 percent (4.83 in the model). A household with a SUD is twice
as likely to develop SMI ($\iota=2.00$), between the hazard ratios of 1.34 and
6.04 in the Danish registers (Hjorth\o j et al., 2018; Nielsen et al., 2017). A
household recovers from SMI with probability 0.066 a year, from the ten year
recovery rate in {\AA}sb{\o} et al.\ (2022).

SMI acts through three channels, each set from its own source. It makes a job
harder to find, with the divisor $\kappa^{ue}_{m}$ calibrated to the SMI share of
the homeless, 25 percent in Gutwinski et al.\ (2021) (24.7 in the model). It
raises the baseline risk of dying 2.50 times (Olfson et al., 2015). And it lowers
the chance of recovering from a SUD, by the multiplier 0.82 (Bradizza et al.,
2006). So the link between mental illness and drug use runs both ways: a SUD
raises the chance of SMI, and SMI makes a SUD more persistent. Untargeted,
households with SMI are 2.64 times as likely as others to use drugs in the model,
against 2.44 in NSDUH 2024.

Two choices are worth stating. SMI does not enter the utility function; it acts
only through these three channels. And housing does not move the SMI chain,
because in the At Home / Chez Soi trial housing had no detectable effect on the
symptoms of homeless adults with SMI (Stergiopoulos et al., 2015).

On your expectation that mental illness weakens Treatment First: in the model it
pulls both ways. Households with SMI use drugs and are homeless at higher rates
than others, as in the data, so Treatment First's drug test falls
disproportionately on them. For those it puts on the street, losing the voucher
raises three year mortality by about 3.6 percentage points, of the same order as
the 4.6 points by which Costantini (2025) finds Housing First lowers the three
year mortality of homeless veterans with mental illness. But a SUD raises the
onset of SMI, so deterring use also prevents some SMI, with the unemployment and
the mortality that come with it. With SMI in the model and calibrated, Treatment
First still raises a newborn's welfare by 2.57 percent. \open The paper does not
report the model without SMI, so I cannot say how much SMI moves that number.

This is in Section 4.7, Block 3 of Section 5.1 (Table 3), and Blocks 2 and 7 of
Section 5.2 (Table 4).
\end{response}

\cmt{1(b)}{Treatment First contains no treatment}
\begin{referee}
``One of the goals of TF as a policy is to encourage addicts to seek treatment for
their addiction (it is even in the name of the policy!). However, this mechanism
is ruled out in the model by the assumption that recovery is random. I found the
argument for ruling out any possibility of this mechanism to be fairly weak given
its centrality to one of the two policies studied.''
\end{referee}
\begin{response}
\concede You are right, and treatment is now in the model. A household with a SUD
chooses each year whether to be in treatment, and treatment doubles its annual
chance of recovering ($g=2.00$, from the remission rates under medication for
opioid use disorder in Paul et al., 2024, and the trials of agonist medication in
Fudala et al., 2003). Treatment costs the household utility, $\bar\tau^{u}$,
calibrated to treatment take up: 24.5 percent of adults with a SUD received
treatment in the past year in NSDUH 2024, against 24.9 in the model. It costs the
government \$9,426 a year, the cost of a residential treatment episode in Roebuck
et al.\ (2003), under both policies, since treatment is publicly funded. Recovery
is also less likely while homeless and with SMI.

Treatment First now requires treatment, as the July 2025 executive order directs:
a household with a SUD keeps its voucher only if it is in treatment and does not
use. In the main experiment the share of households with a SUD who are in
treatment rises from 24.6 to 29.0 percent.

On your expectation that this strengthens Treatment First: in the long run it is
not what drives the result. When households with a SUD cannot keep their voucher
through treatment at all, Treatment First cuts drug use and the share with a SUD
by 29 percent each, homelessness falls by 11 percent, and a newborn gains at
least as much as in the main experiment (footnote 38). Where treatment matters is
the transition. Along the main path homelessness and the share of children in a
shelter never rise, and one of the three reasons Section 6.4 gives is that few
households lose their vouchers, since a household with a SUD keeps its voucher by
entering treatment and not using. When households with a SUD cannot stop using, a
test with no exemption from the first day raises homelessness by up to 14 percent
if it catches every user, and by up to 5 percent if it catches one in five
(Section 6.5).

This is in Sections 4.6 and 4.9, Block 4 of Section 5.1 (Table 3), Section 6.2
(Table 5) and Sections 6.4 and 6.5 (Table 7).
\end{response}

\cmt{1(c)}{The drug price is fixed at one}
\begin{referee}
``However, in the model, the drug price appears to be fixed at one and does not
adjust in the policy counterfactuals. \dots this seems like a strong assumption
that either should be relaxed or better defended.''
\end{referee}
\begin{response}
Relaxed. The drug price is now set in a market that clears against a supply curve
with a constant price elasticity $\varepsilon_s$. The supply shifter is set so
that the price is one in the benchmark and is then held fixed, so the price moves
in any counterfactual that lets the market clear. It clears in the main
experiment and in every period of the transition; other runs hold it at its
benchmark level so that its effect can be read off on its own. The Associate
Editor asks for the same thing, and Referee~3 asks how the price can be varied
when it is normalised to one; this is the answer to both.

Your concern about offsetting is borne out, and the paper reports its size. In
the main experiment the price falls by 5 percent, which draws some households back
into use. With the price held and the chance of drawing a voucher, $\pi_V$,
balancing the budget, Treatment First cuts drug use by 42 percent and a newborn
gains 6.45 percent; with the price clearing as well, use falls by 21 percent and
the gain is 2.57 percent.

The size of the offset depends on $\varepsilon_s$. I know of no direct estimate
for an illicit drug. Cannabis is the closest market I could find: it remains a
Schedule I substance federally, so its suppliers work under federal prohibition,
as an illicit supplier does. I set $\varepsilon_s$ to 1.46, the estimate of Mace
et al.\ (2020) on Washington State administrative data. That is a short run
elasticity, over a month; the long run elasticity of an illicit supply is likely
higher, and studies of illicit drugs take it to be infinite (Becker et al.,
2006). A more elastic supply means a smaller fall in the price and a smaller
offset, so the offset in the paper is more likely too large than too small. The
paper does not report the results at another value of $\varepsilon_s$. On the
demand side, the model's participation elasticity, $-1.32$, is larger in absolute
value than the estimates of about $-0.3$ to $-0.9$, which again works against
Treatment First (2(f)).

This is in Section 4.12, Block 6 of Section 5.1 (footnote 31), Block 5 of Section
5.2, and Section 6.2 (Table 5).
\end{response}

\cmt{1(d)}{Show that the mortality channel is the quantitatively important one}
\begin{referee}
``\dots it seems possible that the labor-market and income-risk consequences of
drug use could also generate a low-income concentration of drug use. The paper
should demonstrate that the claimed mortality/value-of-life channel is actually
the quantitatively important one. A simple way to do this would be to compare the
relationship between income and drug use in the baseline model to an otherwise
similar version of the model in which mortality is not affected by drug use.''
\end{referee}
\begin{response}
\newresult Done in the stylized model of Section 3, which keeps the saving
decision you were concerned about, and the saving decision turns out to be what
decides between the two channels. \open It is not done in the quantitative
model: the paper does not report that model solved again with drug use leaving
mortality unchanged.

In the stylized model a household can save but not borrow, and a SUD raises both
its risk of dying and its risk of unemployment. It is calibrated so that
participation is 8.90 percent and the average user earns 0.74 of the population
average, as in NSDUH 2024. Your comparison is then 0.74 with both channels at
their calibrated strength, against 3.51 with mortality unaffected by a SUD and the
unemployment channel at its calibrated strength. Section 3.3 and Figure 1 run it
one channel at a time, from zero to twice its calibrated strength, with $b$
recalibrated at each point so that participation stays at 8.90 percent and only
who uses changes.

With neither penalty, use concentrates at the top, because the drug is a luxury
good: the average user earns 3.51 times the population average. The mortality
channel alone takes this to 0.64 at its calibrated strength, crossing the average
at 73 percent of that strength, and to 0.27 at twice it. The unemployment channel
alone leaves the average user at 3.51 times the population average, even at twice
its calibrated strength: on its own it does not change who uses at all.

Your conjecture holds in one case, and the paper shows it. With saving switched
off, the right panel of Figure 1, the unemployment channel also reverses the
gradient, to 0.75 at its calibrated strength and 0.27 at twice it. The
difference is self insurance. A household that can save insures itself against
losing its job, so its loss is roughly proportional to its income and costs a
rich household about as much in utility as a poor one. Nothing insures it against
dying, which takes away the whole continuation value of being alive, and that
loss grows with income. Without saving, losing a job takes a household down to
the food assistance floor whatever it earned, so that loss grows with income too.

The result also holds in a version of the stylized model with housing, a
shelter, unemployment insurance and housing vouchers. There, with no consequences
of a SUD, the average user earns 3.54 times the population average; the risk of
dying alone takes this to 0.52 at twice its calibrated strength; the risk of
losing one's job alone leaves it at 3.54; and only with saving switched off does
job loss on its own bring it down, to 0.56 (footnote 10).

In the quantitative model the SUD mortality multiplier, $\kappa_{\phi d}$, is the
parameter calibrated to the income of users, for the reason Section 3 gives, and
Figure 2 compares the model's gradient with NSDUH 2024 (2(c)(i)).

This is in Sections 3.1 to 3.3 (Tables 1 and 2, Figure 1, footnote 10), Block 1
of Section 5.1 (Table 3) and Section 5.2 (Figure 2).
\end{response}

\cmt{1(e)}{The appendix is disorganised and claims more than it establishes}
\begin{referee}
``\dots the appendix is useful for developing intuition, but I found it somewhat
disorganized and narrower than the surrounding discussion suggests. The results
are derived in very specific static settings, while the main model is dynamic \dots
the appendix should be reorganized and made more precise and the main text
discussion should describe it more clearly as intuition rather than as a general
proof of the model's mechanism.''
\end{referee}
\begin{response}
\concede Accepted. The static appendix is gone. In its place, Section 3 is a
stylized model in the body of the paper: two periods of thirty years, a saving
decision, a decision to use, a SUD that raises the risk of dying before the second
period and of unemployment in it, and food assistance as the only transfer. It is
calibrated (Tables 1 and 2) and hits its five targets exactly. The paper says
that its sole purpose is to show why the income gradient in drug use reverses
(1(d)). The welfare results come from the quantitative model of
Section 4, which the paper introduces as the model its numbers come from. The
Associate Editor asked for the same section.

Three footnotes say how far the result travels: it holds with log utility
(footnote 8), when the risk of a SUD rises with use (footnote 9), and in a
version with housing, a shelter, unemployment insurance and vouchers (footnote
10). The only appendix now is Appendix A, which sets out the standard parameters.
The paper does carry the mechanism over to the quantitative model, where it is
the reason the SUD mortality multiplier is calibrated to users' income, but it
does not test the mechanism there (1(d)).

This is in Section 3.
\end{response}

%======================================================================
\section*{2. Calibration and quantification}

\cmt{2(a)}{Model fit never reported}
\begin{referee}
``\dots a notable omission is the model fit. That is, the author never reports the
model's moments even though they would naturally fit in the table. The closest
that the author comes to reporting these moments are by saying, `My benchmark
closely matches all calibration targets.' This is incredibly vague.''
\end{referee}
\begin{response}
\concede You are right, and I agree with 2(b) that without the model's moments
the fit could not be checked. Table 3 now has a Data and a Model
column for each of the 26 parameters that carry the paper's results, and Table
A1 does the same for the 33 standard ones. For every parameter calibrated
internally, 21 in all, the Model column gives the benchmark's
value of its target beside the data. The largest gaps are the income of users, 62
percent of the average in the model against 74 percent in the data (2(c)(i));
renters paying more than half their income, 25.8 against 19.2 percent; and the
unemployment rate, 5.6 against 4.9 percent. The stylized model of Section 3 hits
all five of its targets exactly (Table 2). The sentence you quote is gone.

This is in Section 5.1 (Table 3) and Appendix A (Table A1).
\end{response}

\cmt{2(b)}{Untargeted misses framed too positively}
\begin{referee}
``I find these tables to broadly demonstrating that the model misses substantially
across the most important dimensions presented. \dots most problematic in my eyes
is that the author frames the model as having done a reasonably good job of
matching these moments. \dots Without being able to verify myself \dots I cannot
help but question both the fit of the calibration and also the entire
quantification exercise.''
\end{referee}
\begin{response}
\concede Accepted, and the sentence you quote is gone. Section 5.2 now reports
every untargeted comparison in eight blocks (Table 4), with the model's number
beside the data's. The main misses are these. The homeless in the model are too
concentrated at ages 35 to 44, too few of them are 55 and over, and fewer of them
work or have income (Block 1). Users are younger than in the data, and excess
deaths due to a SUD are a smaller share of deaths at ages 25 to 44 than drug
poisonings are in the data (Block 2). More housed households spend over thirty
percent of their income on rent than renters do in the data, and voucher holders
pay a larger share of their income in rent (Block 3). The bottom half of
households holds a larger share of net worth than in the data (Block 4).
Participation responds more to the drug price than in the estimates (Block 5).
And homelessness passes from parent to child more strongly than in Bassuk et al.\
(1997) (Block 6).

Where the paper can sign what a miss does to the result, it says so. A larger
price response can only work against Treatment First, since the lower drug price
under Treatment First draws more households back into use (Block 5). With
transmission at the low end of the evidence on the SUD and above it on
homelessness, the model loads the externality on children towards Housing First
more than Treatment First (Block 6). For the other misses the paper does not sign
the effect. With the fit of every target now in Tables 3 and A1 (2(a)), you can
check both the targeted and the untargeted fit yourself.

This is in Section 5.2 (Table 4).
\end{response}

\cmt{2(c)(i)}{Homeless and users are too poor and too detached from work}
\begin{referee}
``In the model, all homeless individuals are unemployed, compared with 40 to 52
percent in the data, and average income among the homeless is essentially zero,
compared with 19.1 percent of average income in the data. Similarly, the average
income of drug users is 41 percent of average income in the model, compared with
74 percent in the data. \dots a model in which these groups are much poorer and
less attached to employment than in the data may evaluate HF and TF at the wrong
margins.''
\end{referee}
\begin{response}
\concede You are right, and the paper now targets the 74 percent you cite. In
both models a user is a household whose use of an illicit drug amounts to a
disorder, mild included, the counterpart of adults aged 26 and over with a drug
use disorder of any severity in NSDUH 2024, whose average income is 74 percent of
the population average. A household with a SUD is one whose disorder is moderate
or severe, and use that never reaches a disorder is not modelled (footnote 26).
The stylized model of Section 3 matches the 74 percent exactly. In the
quantitative model the SUD mortality multiplier, $\kappa_{\phi d}$, is calibrated
to it, and the model reaches 62 percent (Table 3). Users in the model are still
poorer than in the data, but by much less than in the version you read.

The calibration targets only the average income of users. Figure 2, new in
Section 5.2, compares the whole distribution of use by income with NSDUH 2024 and
shows where the remaining gap comes from. As in the data, use is highest at the
bottom and lowest at the top: 18.3 percent of households with less than 0.22 of
the mean income use drugs, against 13.7 percent in the data, and 4.2 percent of
those with more than 1.62 times the mean, against 4.8 percent. Between 0.43 and
1.62 times the mean, use in the model is flatter and lower than in the data, 5.9
to 6.3 percent against 7.2 to 10.2. Both differences lower the average income of
users in the model.

The homeless are also less deprived than in the version you read, though still
more than in the data. Of the homeless in the model, 21.4 percent are employed,
against 46 percent in Meyer et al.\ (2024), and their income is 6.5 percent of the
average, against 19.1 percent. Both are untargeted.

On the margins you list, Section 5.2 now says why the income of users matters for
the experiments: the vouchers are means tested, so the lower in the income
distribution users sit, the more of them are eligible for a voucher, and the more
of them, and of their children, Treatment First's test reaches. The model no
longer has a labour supply choice, so employment works through the chances of
losing and finding a job, which depend on productivity, homelessness, a SUD and
SMI; untargeted, unemployment among households with a SUD is 14.5 percent in the
model against 16.2 percent in NSDUH 2024. And the rise in tax revenue under
Treatment First comes mainly from children who start adult life with higher
productivity, and only to a lesser extent from fewer adults out of work because of
a SUD.

\open The misses remain, and the paper does not show how the results would change
if users and the homeless were as well off as in the data.

This is in Block 1 of Section 5.1 (Table 3), Blocks 1 and 2 of Section 5.2 (Table
4, Figure 2), Sections 4.3 and 4.8, and Section 6.2.
\end{response}

\cmt{2(c)(ii)}{The model's homeless are too young}
\begin{referee}
``In the model, 73.1 percent of homeless individuals are between ages 26 and 44,
compared with 48.1 percent in the California data. Conversely, only 9.3 percent of
the model's homeless population is age 55 or older, compared with 27.1 percent in
the data. \dots If the model shifts homelessness toward younger ages, it may
overstate the long-run labor-market and addiction-response channels while
understating the role of policies for older homeless individuals.''
\end{referee}
\begin{response}
\concede You are right, and the miss has narrowed but not closed. The age
distribution of the homeless is now compared with the national point in time
count of the Department of Housing and Urban Development for 2024. The
youngest band, 25 to 34, is the target of the endowment at entry to adulthood,
$T_0$ (26.4 percent in the model against 26.0 in the data); the other four bands
are untargeted. The model still has too many of the homeless at 35 to 44, 37.9
percent against 27.2, and too few at 55 and over, 11.5 percent against 25.9.
Users, too, are younger in the model than in the data (S1).

Section 5.2 now says why the age of users matters for the experiments: both
policies reach a child through its parent's decisions, so the age of users
decides how many children their use touches. The paper's check on how many children the
externalities reach is the share of children living with a parent who has a SUD,
5.8 percent in the model against 4.6 percent in NSDUH 2024 (Block 6). \open The
paper does not show how the results would change with an older homeless
population, and I do not claim the gap is harmless.

This is in Block 2 of Section 5.1 (Table 3) and Blocks 1, 2 and 6 of Section 5.2
(Table 4).
\end{response}

\cmt{2(d)}{Match the policy relevant elasticities}
\begin{referee}
``Relevant objects include the response of drug use to voucher generosity or
conditionality, the response of homelessness to voucher access, the response of
labor supply and employment to vouchers, and the response of recovery or addiction
persistence to housing status. \dots the RCT literature cited in the introduction
may provide identified elasticities that the model can match.''
\end{referee}
\begin{response}
Of the responses you list, two are now matched or compared with an estimate: the
response of recovery to housing status, which is a calibration target, and the
response of drug use to conditionality, through two estimates. The paper adds a
third, the response of mortality to losing a voucher. The other three are not
compared with an estimate.

\begin{itemize}
\item \emph{Recovery and housing status.} $\mu_H$, the multiplier on the chance
of recovering from a SUD while homeless, is calibrated to Milby et al.\ (2005),
who find that housing improves the chance of recovering by 47.3 percent (1.47 in
the data, 1.45 in the model).
\item \emph{Drug use and conditionality.} Milby et al.\ (2005) also randomise
housing conditional on abstinence against housing without a condition, among
homeless adults with a cocaine disorder in treatment: 8.32 against 6.60
consecutive weeks of abstinence, a ratio of 1.26. The model's counterpart, the
recovery hazard of voucher holders with a SUD under the condition against that of
housed households with a SUD and no voucher, in the first year, is 1.40,
untargeted.
\item \emph{Deterrence of a sobriety condition.} The detection rate of Treatment
First's drug test is set so that the test deters use as much as in the closest
evidence, the US military's random testing with dismissal, after which service
members who stayed were 61 percent less likely than civilians to have used
marijuana in the past year; counting those who left, the fall was 11 to 45
percent (Mehay and Pacula, 1999). In the model, 62 percent of the users who face
the test stop when it is introduced (4(a)).
\item \emph{Housing and mortality.} For households with SMI whom Treatment First's
test puts on the street, losing the voucher raises three year mortality by about
3.6 percentage points in the model, with the drug price and $\pi_V$ held, against
the 4.6 points by which enrolling in Housing First lowers the three year mortality
of homeless veterans with mental illness in Costantini (2025), untargeted. The
paper compares the model with his Housing First estimate only; his finding that
the Treatment First programmes of the Department of Veterans Affairs have no
lasting effect on mortality is cited in Section 2 but not compared with the
model.
\item \emph{Labour supply and employment.} The model has no labour supply choice:
earlier versions had one, and it had a negligible effect on the results
(footnote 19). Housing affects employment through $\kappa_H$, set at 0.20 from
Desmond and Gershenson (2016): in a shelter, separations are $1+\kappa_H$ times
as frequent and job finding $1-\kappa_H$ times as fast.
\end{itemize}

\open The paper does not compare the model's response of drug use to voucher
generosity, of homelessness to voucher access, or of employment to vouchers with
an estimate, although the Family Options Study it cites in Section 2 measures the
effect of an unconditional subsidy on later homelessness (Gubits et al., 2016,
2018).

This is in Blocks 2 and 4 of Section 5.1 (Table 3), Blocks 7 and 8 of Section 5.2
(Table 4), the paragraph on the screen at the start of Section 6, and Section 4.4.
\end{response}

\cmt{2(e)}{The price elasticity moment is not replicable}
\begin{referee}
``\dots I found the description of the calculation of this moment to be extremely
vague, to the extent that I would not be able to replicate it. This should be
clarified. I would appreciate an appendix describing the model definitions of the
moments more thoroughly in general.''
\end{referee}
\begin{response}
The participation elasticity is no longer a calibration target, so the
calibration does not depend on how it is computed. The model's participation and
intensive elasticities are reported as untargeted moments, $-1.32$ and $-0.63$,
against about $-0.3$ to $-0.9$ and $-0.3$ to $-0.8$ in the literature. The
response of drug use to Treatment First's test itself, which the paper keeps
separate from the response to the price, is matched instead to the military's
random testing (2(d), 4(a)).

On the general request: Tables 3, 4 and A1 give, for each parameter and each
check, the moment and its source, and the notes to Table 4 and to Figure 2 define
the model's counterpart where it is not a simple share. The replication package
carries every dataset and a data appendix, \pname{data/DATA\_APPENDIX.md}, which
documents the dataset, year, universe and computation behind every data estimate.
\open The paper has no appendix that defines the model's moments, and it does not
spell out how the model's two price elasticities are computed.

This is in Block 5 of Section 5.2 (Table 4) and the opening of Section 5.
\end{response}

\cmt{2(f)}{Show identification numerically, and show robustness to the elasticity range}
\begin{referee}
``I did find the identification explanation in Section 4.2 to be quite helpful.
\dots I would like this to also be shown numerically. This could be done in many
ways: reporting the Jacobian of the moment vector with respect to the parameter
vector, reporting the sensitivity statistic from Andrews, Gentzkow, and Shapiro
(2017), through a figure showing how perturbing parameters locally around the
baseline calibration affects the associated moments. \dots take the elasticity of
drug use with respect to price which has a range of plausible estimates from
$-0.36$ to $-0.9$ \dots whether or not the results are robust to this choice \dots
should be demonstrated.''
\end{referee}
\begin{response}
\open The paper does not report a Jacobian, the statistic of Andrews, Gentzkow and
Shapiro (2017), or a perturbation figure. What it does is narrower. Each
parameter's paragraph in Section 5.1 says what the parameter controls, what
happens as it rises and which moment it is matched to, in that order, and Table 3
now reports the model's value of each moment beside the data (2(a)).

On the price elasticity: you are right that its range is wide and that it matters
for the policy result. It is no longer a target (2(e)), so the calibration does
not depend on where in the range the target is placed. The model's participation
elasticity, $-1.32$, is above the range in absolute value, and the paper argues
that this is conservative: a larger price response can only work against
Treatment First, since the lower drug price under Treatment First draws more
households back into use. The paper does not solve the model with an elasticity
inside the range.

For the other inputs that are measured imprecisely, the robustness checks the
paper reports are listed under 4(d).

This is in Sections 5.1 and 5.2 (Tables 3 and 4).
\end{response}

%======================================================================
\section*{3. Policy analysis}

\cmt{3(a)}{Loss of consortium is not integrated into the model}
\begin{referee}
``If premature deaths impose welfare losses on other individuals, then these
losses should in principle affect the value of life and agents' decisions (such as
the decision to use drugs). Instead, the paper adds the loss-of-consortium effect
only after solving the model. I find this to be an internal inconsistency.''
\end{referee}
\begin{response}
\concede You are right, and the adjustment is gone: nothing is added to welfare
after the model is solved. I did not build a loss of consortium into the model; I
dropped it. Referee~2 pointed to two effects on others that the previous version
missed, a disamenity from homelessness and parents' concern for their children.
The paper builds a version of the second, inside the household's problem.

Each household has a child, who lives with it for its first 25 years, from the
parent's age 26 to 50, shares its consumption and dwelling but not the drug, and
starts adult life with lower productivity if the parent had a SUD while the child
lived at home. The parent values its child at the weight $\omega_c$, both in the
child's flow utility during childhood (equation 12) and through what the parent's
SUD costs the child's value as an adult, charged when the child leaves home
(equations 24 and 25). So the child's welfare enters the parent's choices,
including the choice to use, as you argue a loss to others should; only the share
$1-\omega_c$ that the parent leaves out is an externality. The welfare criterion
counts the child's childhood at its own weight, and Section 6.3 measures the part
of the gain the parent does not count (Table 6).

The disamenity is not in the benchmark. I solved the model with one, at the value
in Loubi\`ere et al.\ (2020), about 0.2 percent of average earnings a year, and no
long run gain in Table 5 moves by more than 0.14 percentage points.

This is in Sections 4.2, 4.8 and 4.10, Sections 6.1 and 6.3 (Table 6), and
footnote 4.
\end{response}

\cmt{3(b)}{The adjustment is ad hoc and a truly conservative value is zero}
\begin{referee}
``$\kappa_{cons}$ was plucked out of thin air, yet makes an enormous impact on the
final policy rankings. To play devil's advocate on the point that the chosen value
was conservative, a truly conservative value for a parameter plucked out of thin
air such as this would be zero, which results in a different policy
recommendation.''
\end{referee}
\begin{response}
\concede You are right, and the parameter is gone. The paper's first experiment
is the case you describe, with no cost of a household's drug use to anyone else
beyond the budget. The children are removed from welfare and from the parent's
preferences, and the child's productivity no longer depends on what its parent
did, so a household's drug use can harm only itself and the budget, and only
moral hazard is left to justify a condition. In that setting Treatment First
loses: a newborn is worse off by 1.17 percent of lifetime consumption with the
drug price and $\pi_V$ held, by 1.09 percent when $\pi_V$ rebalances the budget,
and by 0.74 percent when the drug price clears as well. The paper states this as a result:
moral hazard alone does not justify conditionality.

The case for Treatment First then rests on the externalities on children, and the
two parameters that set them are each tied to evidence. The child's earnings
penalty, $\delta_d=0.24$, is the estimate of Nilsson (2017), to my knowledge the
only natural experiment on how a parent's substance shock affects a child's adult
earnings. It measures exposure to alcohol in the womb, not a parent's disorder
after birth, which no study measures directly. As an indirect check, the model's
odds that a child of a parent with a SUD acquires one, 1.46, are at the low end
of the 1.44 to 2.51 in the literature, which suggests that the penalty, if
anything, is at the low end as well. The parent's weight on its child,
$\omega_c=0.085$, is set by behaviour: in NSDUH 2024, unmarried adults living with
one other adult and no child are 17.5 percent more likely to have used illicit
drugs in the past year than lone parents with one child, reweighted to the same
age, sex, marital status and family income, and $\omega_c$ is the weight at which
removing a parent's concern for its child raises use in the model by as much
(17.7 percent).

The paper reports that the sign of the result survives a change in either. With
$\omega_c$ at 0.69, the upper end of the literature (Nishiyama, 2000), a newborn
gains 3.18 percent, though only 23 percent of that gain is then an externality;
with the earnings penalty set to zero, Treatment First still raises welfare; and
the earnings externality is 18 percent of the gain in the main experiment.

This is in Section 2, Section 6.2 (Table 5), Block 5 of Section 5.1 (Table 3),
Block 6 of Section 5.2 (Table 4), and footnotes 5, 28 and 30.
\end{response}

\cmt{3(c)}{Tone down the language}
\begin{referee}
``Universal TF is not unambiguously preferred under all welfare criteria
considered in the paper. \dots An example of language that should be softened is
the conclusion's statement that the debate is resolved in favor of universal TF.''
\end{referee}
\begin{response}
\concede Accepted. The paper no longer compares universal programmes, and it does
not claim the debate is resolved. Every experiment keeps the programme rationed,
as it is, and the answer is stated conditionally: Section 7 says that housing
assistance should require sobriety if the externalities on children are counted,
and that moral hazard alone is not a reason for it.

The paper also says who loses. As the trials imply, most households with a SUD
are worse off under Treatment First: 74 percent of those who have one when the
reform lands, with the drug price and $\pi_V$ held as in a trial, and 67 percent
when both clear. And the path to the long run depends on how the reform is
designed (3(d)).

This is in Section 7, Section 6.4 and Section 6.5 (Table 7).
\end{response}

\cmt{3(d)}{Solve proper transition paths}
\begin{referee}
``The author should solve for proper transition paths. \dots I find the current
strategy, of a several-period PE adjustment, to be a half-measure that is limited
in its informativeness and opens up more questions than it answers.''
\end{referee}
\begin{response}
\newresult Done. The quantitative model's transition from Housing First to
Treatment First is solved and reported. The reform lands as an unexpected
permanent change at the start of year 1, the path starts from the benchmark, and
the drug price and $\pi_V$ clear the drug market and the budget in every period,
for sixty years, to the Treatment First steady state.

You expected the slow adjustment of the SUD to matter, and it does, though not in
a way that reverses the result. Drug use falls on impact, 81 percent of its long
run fall in the first year. The share with a SUD falls only as fast as recovery
allows, half of its long run fall by year 11 and nine tenths by year 27, and the
drug price stays above its long run level for 33 years. Homelessness and the share
of children in a shelter never rise: homelessness falls from 0.62 to 0.60 percent
of households by year 10, and the share of children in a shelter from 1.03 to 1.00
percent. Section 6.4 gives three reasons. Few households lose their vouchers: at
most 0.13 percent of households are excluded from the programme in any year,
because a household with a SUD keeps its voucher by entering treatment and not
using, and the test catches a user with a probability of only 20 percent a year.
Drug use falls from the first year. And after the first three years $\pi_V$ stays
above its Housing First level. Children born at the reform gain 1.45 percent of
lifetime consumption, less than the 2.57 percent of those born in the long run,
because they grow up while more parents still have a SUD and before $\pi_V$ has
fully risen. Of the adults alive at the reform, 11 percent are worse off.

The design of the reform matters for the first years, and Section 6.5 shows how.
When households with a SUD cannot stop using, a test with no exemption from the
first day that catches every user raises homelessness by up to 14 percent and the
share of children in a shelter by up to 11 percent, and they are above their
Housing First levels in periods 1 to 15 and 2 to 11. If the test catches one in
five users, the rises are up to 5 and 3 percent. Announcing the test five periods
ahead lowers the rise but does not remove it. Exempting households in treatment
for the first five periods keeps the share of children in a shelter from ever
rising above its Housing First level, and holds the rise in homelessness to 1.6
percent, when the test catches one user in five. Children born at the reform gain
in every one of these alternatives.

This is in Section 6.4 (Figure 3) and Section 6.5 (Table 7).
\end{response}

%======================================================================
\section*{4. Sensitivity}

\cmt{4(a)}{Imperfect detection}
\begin{referee}
``TF assumes that drug use is perfectly observed and perfectly enforceable. \dots
If drug use were imperfectly observed, this would likely reduce the effectiveness
of TF. The paper should analyze the sensitivity of the results to imperfect
detection.''
\end{referee}
\begin{response}
Added. Under Treatment First every income eligible voucher holder, and every
income eligible household that draws a voucher, is tested, and one that fails the
condition is
caught with probability $p_{\mathrm{det}}$; with probability $1-p_{\mathrm{det}}$
it keeps its voucher. A household that is caught loses the voucher from the next
period and is excluded from the programme for three years in expectation, the
period for which the Housing Choice Voucher programme already bars a household
evicted for drug related criminal activity.

$p_{\mathrm{det}}$ sets how strongly the test deters, so I set it to the closest
evidence, the military's random testing (2(d)). At $p_{\mathrm{det}}=20$ percent,
62 percent of the users who face the test stop when it is introduced, against the
61 percent of Mehay and Pacula (1999) among service members who stayed.
Enforcement is far from perfect: the test catches only 20 percent of users a year.

The sign of the result survives weaker detection. At half the rate, 10 percent, 41
percent of the users who face the test stop, inside the range of 11 to 45 percent
that Mehay and Pacula find once those who left the military are counted. In the
long run Treatment First then cuts drug use by 14 percent and homelessness by 6
percent, against 21 and 11 percent at the calibrated rate, and a newborn gains
1.65 percent, against 2.57. At the other end, Section 6.5 compares catching every
user with catching 20 percent, along transitions in which households with a SUD
cannot stop using (Table 7).

This is in Section 4.9 and the paragraph on the screen at the start of Section 6
(footnotes 33 and 35).
\end{response}

\cmt{4(b)}{False positives}
\begin{referee}
``\dots the paper should consider the consequences of false positives. If some
non-users lose voucher access because of erroneous tests, then TF may impose
substantial costs on individuals who are not violating the sobriety requirement.''
\end{referee}
\begin{response}
\open Not added. In the model a household that does not use, and is in treatment
if it has a SUD, is never flagged, because a positive test is confirmed in a
laboratory before it has consequences. That is the paper's reason for leaving
false positives out. It does not cover the other ways a household that complies
could lose its voucher, such as administrative error, and the paper does not
price those.

This is in the paragraph on the screen at the start of Section 6.
\end{response}

\cmt{4(c)}{Compliance costs, stigma, administrative burden}
\begin{referee}
``I found the argument that drug tests would cost only billions of dollars per
year to implement to both be most likely a massive underestimation that ignored
administrative overhead and assumed far more efficiency than is warranted of
government policy, and also to be too large a number to entirely ignore in the
first place.''
\end{referee}
\begin{response}
You are right that the cost should not be ignored, and it is no longer outside
the model. Each test costs the government \$35, 0.031 percent of average earnings, the figure in
the Texas Legislative Budget Board's fiscal note for Senate Bill 11, which is the
upper end of the estimates, and it is paid for every test (4(a)). When $\pi_V$
adjusts to hold the government's net revenue at its benchmark level, as in the
main experiment, every dollar spent on testing is a dollar less for vouchers.

Treatment, which Treatment First requires of households with a SUD, is a second
cost. The government pays \$9,426 a year for each household in treatment, under
both policies, and as take up rises under Treatment First, part of what the test
saves goes on treatment rather than on vouchers. The household bears a cost too:
treatment costs it utility, and that cost is calibrated to observed take up.

\open What the paper does not do: it has no stigma, no administrative cost beyond
the cost per test, and no sensitivity run on the cost per test.

This is in Sections 4.9 and 4.11, Block 4 of Section 5.1, and the paragraph on
the screen at the start of Section 6.
\end{response}

\cmt{4(d)}{Sensitivity to hard to identify parameters}
\begin{referee}
``The paper should report sensitivity to parameters that are critical for the
policy analysis yet are difficult to identify such as the parameters governing
addiction transitions, recovery probabilities, addiction-related mortality,
homelessness-related mortality, and the employment consequences of addiction and
homelessness.''
\end{referee}
\begin{response}
\open The paper does not report the results with any of these parameters changed.
What it reports for each is where the value comes from and how the model compares
with independent evidence.

\begin{itemize}
\item \emph{Addiction transitions.} $\pi_d$ is calibrated to the share of users
with a SUD, 41.3 percent in NSDUH 2024 (42.6 in the model). In the stylized model
the results hold when the risk rises with use (footnote 9), and in earlier
versions of the quantitative model letting it rise with use had little impact on
the results (footnote 20).
\item \emph{Recovery.} The baseline rate, 0.025 a year, is the hazard implied by
the ten year remission rate in Vaillant (1995); the effect of housing, $\mu_H$, is
calibrated to Milby et al.\ (2005); and the effect of treatment, $g=2.00$, is
from Paul et al.\ (2024) and Fudala et al.\ (2003). The model's response of
recovery to a sobriety condition, 1.40, is close to Milby's 1.26, untargeted
(2(d)). And the long run result does not depend on treatment: when households
with a SUD cannot respond at all, a newborn gains at least as much (footnote 38).
\item \emph{SUD mortality.} $\kappa_{\phi d}=4.00$, calibrated to the income of
users, lies between the mortality rate ratios of 3.91 in Plana-Ripoll et al.\
(2019) and 6.6 in Heiberg et al.\ (2018). Untargeted, excess deaths due to a SUD
are 14.5 percent of deaths at ages 25 to 44 in the model, against 20.0 percent
from drug poisonings in the Centers for Disease Control and Prevention (CDC)
WONDER database.
\item \emph{Homelessness mortality.} $\kappa_{\phi H}$ multiplies baseline
mortality by 1.60 at every age, the ratio for homeless adults against poor housed
adults in Meyer et al.\ (2023). Untargeted, for households with SMI whom the test
puts on the street, losing the voucher raises three year mortality by about 3.6
percentage points, against Costantini's 4.6 (2(d)).
\item \emph{Employment.} $\kappa^{ue}_{d}$ is calibrated to the SUD share of the
homeless, 21.7 percent in Gutwinski et al.\ (2021) (20.2 in the model), and
unemployment given a SUD is 14.5 percent in the model against 16.2 in NSDUH 2024,
untargeted. The employment penalty of homelessness, $\kappa_H=0.20$, is from
Desmond and Gershenson (2016).
\end{itemize}

The robustness checks the paper does report are on the detection rate (halved,
footnote 35), the parent's weight on its child ($\omega_c=0.69$, footnote 28),
the child's earnings penalty (set to zero, footnote 5), a disamenity from
homelessness (footnote 4), and whether households with a SUD can respond to the
condition at all (footnote 38 and Section 6.5). In each, Treatment First still
raises a newborn's welfare.

This is in Section 5.1 (Table 3), Section 5.2 (Table 4) and the footnotes named
above.
\end{response}

%======================================================================
\section*{Smaller comments and typos}

\cmt{S1}{Why a maximum experimentation age?}
\begin{referee}
``I do not understand why there is assumed to be an age at which drug
experimentation will not occur. This seems inconsistent with Table 4(a), which
shows U-shaped drug use by age.''
\end{referee}
\begin{response}
\open The limit stays in the model, and the paper explains why. It is at model
age 55, age 80 in a life that runs to 85, and it stops only households without a
SUD from starting to use; a household that has a SUD at that age can keep using.
So it limits when use can begin, not how long it lasts, and it binds only after
age 80. Without it, using at old ages has little or no consequence, because too
little life remains for the disorder to damage, and the model produces too much
use in old age. The cause is that death is certain at the last age in the model,
so close to it a household knows how little life remains, whereas in real life
most people do not know when they will die.

The age profile of users is now reported against NSDUH 2023, and the model misses
it at both ends. Users in the model are younger than in the data: 67 percent of
them are 26 to 34, against 45 percent, and their mean age is 34.6 years, against
39.6. At 55 to 64 and at 65 and over the model has 5 and 2 percent of users,
against 12 and 10 percent in the data.

This is in Section 4.6 (footnote 21), Block 13 of Appendix A (Table A1) and Block
2 of Section 5.2 (Table 4).
\end{response}

\cmt{S2}{Tables should show levels and percentage changes}
\begin{referee}
``The results tables would be easier to interpret if they included both percentage
changes and baseline/counterfactual levels for the main policy outcomes.''
\end{referee}
\begin{response}
Done in part. Table 5 reports, for every outcome, its level under Housing First
and under each version of Treatment First side by side, with and without the
externalities, and Section 6.2 gives the percentage changes in the text. Table 7
reports the transitions as changes against Housing First: the largest percent rise
in homelessness and in the share of children in a shelter, and the periods in
which each is above its Housing First level. \open Table 5 does not have a
separate column of percentage changes.

This is in Section 6.2 (Table 5) and Section 6.5 (Table 7).
\end{response}

\cmt{S3 to S10}{Typographical corrections}
\begin{referee}
``well-aware'' should be ``well aware''; equation (9) should end with a comma;
remove the comma after the citation in calibration target 19; footnote 7 ``our
calibration'' should be singular; the Abramson citation lists forthcoming twice;
``Kinneruda'' and ``Kragh-S\o rensena'' carry spurious trailing letters; Goering et
al.\ has ``Guido. Powell ,''; Degenhardt et al.\ has ``Hamamd Ali'' for
``Hammad''.
\end{referee}
\begin{response}
The paper has been rewritten, and none of these passages remains: the sentence
with ``well aware'', the old equation (9), the old calibration target 19 and the old
footnote 7 are gone, and the paper is in the first person singular throughout. In
the references, Abramson (2026) is listed once, as published online in the
\emph{Journal of Finance}, Goering et al.\ (2014) no longer carries the error, and
Karlman et al.\ and Degenhardt et al.\ are no longer cited.

This is in the references.
\end{response}

\end{document}
